\section{Theory}
\label{app:theory}

In this section, we present the proof of \cref{thm:ordering_sensitivity} and an example showing that the difference between the two decompositions can be made arbitrarily large even when the mean gap is fixed.

\subsection{Proof of \cref{thm:ordering_sensitivity}}
\label{sec:sign_flip_proof}
\begin{proof}
Using \cref{def:unexplained_sign_flip}, we want to show that 
\[
\sign\!\left( \mu_K^T \Delta\beta + \Delta\alpha \right)
\neq
\sign\!\left( \mu_H^T \Delta\beta + \Delta\alpha \right)
\iff
\begin{aligned}
&\mu_H^T \Delta\beta \neq \mu_K^T \Delta\beta, \\
&\min\{\mu_H^T \Delta\beta, \mu_K^T \Delta\beta\}
<
-\Delta\alpha
<
\max\{\mu_H^T \Delta\beta, \mu_K^T \Delta\beta\}.
\end{aligned}
\]


$\Rightarrow$: Suppose that $\sign\left( \mu_K^T \Delta\beta + \Delta\alpha \right) \neq \sign\left( \mu_H^T \Delta\beta + \Delta\alpha \right)$. If $\mu_H^T \Delta\beta = \mu_K^T \Delta\beta$, then both $\mu_g^T \Delta\beta + \Delta\alpha$ would have the same sign, which contradicts our assumption of opposite signs. Thus, we must have $\mu_H^T \Delta\beta \neq \mu_K^T \Delta\beta$.

Next, suppose that $\mu_K^T \Delta\beta + \Delta\alpha >0$ and $\mu_H^T \Delta\beta + \Delta\alpha <0$. Then, 
\[
\mu_K^T \Delta\beta > -\Delta\alpha \quad \text{and} \quad \mu_H^T \Delta\beta < -\Delta\alpha.
\]
Combining these two inequalities, we have 
\begin{equation}
  \label{eq:intermediate_ineq1}
  \mu_H^T \Delta\beta < -\Delta\alpha < \mu_K^T \Delta\beta.
\end{equation}
Similarly, if $\mu_K^T \Delta\beta + \Delta\alpha <0$ and $\mu_H^T \Delta\beta + \Delta\alpha >0$, then
\begin{equation}
  \label{eq:intermediate_ineq2}
  \mu_K^T \Delta\beta < -\Delta\alpha < \mu_H^T \Delta\beta.
\end{equation}
Combining \cref{eq:intermediate_ineq1,eq:intermediate_ineq2}, we conclude that
\[
  \min\{\mu_H^T \Delta\beta, \mu_K^T \Delta\beta\} < -\Delta\alpha < \max\{\mu_H^T \Delta\beta, \mu_K^T \Delta\beta\}.
\]

$\Leftarrow$; We now show the converse direction. Suppose that:

$(1)$ $\mu_H^T \Delta\beta \neq \mu_K^T \Delta\beta$ and 

$(2)$ $\min\{\mu_H^T \Delta\beta, \mu_K^T \Delta\beta\} < -\Delta\alpha < \max\{\mu_H^T \Delta\beta, \mu_K^T \Delta\beta\}$. 

We want to show that $\sign\left( \mu_K^T \Delta\beta + \Delta\alpha \right) \neq \sign\left( \mu_H^T \Delta\beta + \Delta\alpha \right)$.
$(2)$ implies that one of the following must be true:
\[
  \mu_H^T \Delta\beta < -\Delta\alpha < \mu_K^T \Delta\beta, \quad \text{or} \quad \mu_K^T \Delta\beta < -\Delta\alpha < \mu_H^T \Delta\beta.
\]
If the first one holds, then $\mu_H^T \Delta\beta + \Delta\alpha <0$ and $\mu_K^T \Delta\beta + \Delta\alpha >0$, that is, the unexplained components have opposite signs. If the second holds, a symmetric argument yields the same conclusion as desired. 

% First, consider the case where $\sign \mu_K^T \Delta\beta = \sign \mu_H^T \Delta\beta$.
% In this case, the addition of $\Delta\alpha$ must be small enough to flip one of the signs of the $\mu_g^T \Delta\beta$ but so large that it flips the other.
% As we have assumed $\mu_K^T \Delta\beta \neq \mu_H^T \Delta\beta$, one of the $\mu_g^T \Delta\beta$ will be lesser in magnitude; without loss of generality, suppose this is $g = H$.
% That is, $\abs{\mu_H^T \Delta\beta} < \abs{\mu_K^T \Delta\beta}$.
% Now, we have also assumed that $\sign \Delta\alpha \neq \sign \mu_K^T \Delta\beta$ and $\abs{\mu_H^T \Delta\beta} < \abs{\Delta\alpha} < \abs{\mu_K^T \Delta\beta}$.
% We conclude that $\sign( \mu_H^T \Delta\beta + \Delta\alpha) \neq \sign( \mu_K^T \Delta\beta + \Delta\alpha)$.

% Next, consider the case where $\sign \mu_K^T \Delta\beta \neq \sign \mu_H^T \Delta\beta$.
% In this case, the addition of $\Delta\alpha$ to the $\mu_g^T \Delta\beta$ must be small enough such that it does not change the sign of either one.
% Without loss of generality, suppose that $\sign \Delta\alpha \neq \sign \mu_H^T \Delta\beta$.
% Then the assumed condition $\abs{\mu_H^T \Delta\beta} > \abs{\Delta\alpha}$ is enough to ensure that $\sign(\mu_H^T \Delta\beta) = \sign(\mu_H^T \Delta\beta + \Delta\alpha) \neq \sign(\mu_K^T \Delta\beta + \Delta\alpha)$.
\end{proof}

% ---------------------------------------------------------
% \mqedit{
% \subsection{Decision tree for sign flips}
% \label{sec:decision_tree}

% To shed more light on \cref{thm:ordering_sensitivity}, we disentangle the condition in \cref{eq:sign_flip_condition} into a decision tree of when a sign flip can occur: 

% \begin{algorithm}[H]
%     \caption{Unexplained component sign-flip}
%     \label{alg:decision_tree}
% \begin{algorithmic}  
%     \If{$\mu_H^T \Delta\beta \neq \mu_K^T \Delta\beta$}
%         \If{$\sign(\mu_H^T \Delta\beta) = \sign(\mu_K^T \Delta\beta)$}
%             \If{$\sign(\Delta\alpha) \neq \sign(\mu_H^T \Delta\beta)$}
%                 \If{$|\mu_H^T \Delta\beta| < |\Delta\alpha| < |\mu_K^T \Delta\beta|$ \textbf{or} $|\mu_K^T \Delta\beta| < |\Delta\alpha| < |\mu_H^T \Delta\beta|$}
%                     \State A sign flip occurs.
%                 \EndIf
%             \EndIf
%         \ElsIf{$\sign(\mu_H^T \Delta\beta) \neq \sign(\mu_K^T \Delta\beta)$}
%             \If{$\sign(\Delta\alpha) = \sign(\mu_H^T \Delta\beta)$ \textbf{and} $|\mu_K^T \Delta\beta| > |\Delta\alpha|$}
%                 \State A sign flip occurs.
%             \ElsIf{$\sign(\Delta\alpha) \neq \sign(\mu_H^T \Delta\beta)$ \textbf{and} $|\mu_H^T \Delta\beta| > |\Delta\alpha|$}
%                 \State A sign flip occurs.
%             \EndIf
%         \EndIf
%     \Else
%         \State There is no sign flip.
%     \EndIf
% \end{algorithmic}
% \end{algorithm}

% }


\subsection{Unbounded sensitivity at fixed mean gap}
 \label{ex:unbounded}
In this section we present a toy example where we can have arbitrarily large flip magnitudes. Consider the following example in dimension \( 1 \), where we assume we have a parameter \( L > 0 \) which controls the magnitude of the covariate shift, and \( \varepsilon > 0 \) which controls the deviation in slope coefficients. We assume the mean vector is given by \( (\mu_K, \mu_H) = (0, L) \), and the slope parameters deviate symmetrically from \( \frac{1}{L} \), so that \( \beta_H = \frac{1}{L} + \frac{\varepsilon}{2} \), \( \beta_K = \frac{1}{L} - \frac{\varepsilon}{2} \). Then, the explained and unexplained components under both reference groups are:
\begin{align*}
    \EHOB &= \Delta\mu \cdot \beta_H = L\left(\frac{1}{L} + \frac{\varepsilon}{2}\right) = 1 + \frac{L\varepsilon}{2}, \quad \UH = \mu_K \cdot \Delta\beta = -\frac{L\varepsilon}{2}, \\
    \EKOB &= \Delta\mu \cdot \beta_K = L\left(\frac{1}{L} - \frac{\varepsilon}{2}\right) = 1 - \frac{L\varepsilon}{2}, \quad \UK = \mu_H \cdot \Delta\beta = +\frac{L\varepsilon}{2}.
\end{align*}

Hence, the difference between the two decompositions is \( \EHOB - \EKOB = L\varepsilon \), which can be made arbitrarily large as \( L \to \infty \), even though the total mean gap is always fixed at \( \Delta_Y = 1 \). Similarly, the unexplained components differ by \( \UK - \UH = L\varepsilon \).

\begin{center}
\begin{tabular}{ccccccc}
\toprule
\(L\) & \(\varepsilon\) 
& \(\EHOB\) & \(\EKOB\) 
& \(\UH\) & \(\UK\) 
& Flip? \\
\midrule
20    & 0.1  & \(2\)     & \(0\)      & \(-1\)    & \(1\)     & \checkmark \\
200   & 0.1  & \(11\)    & \(-9\)     & \(-10\)   & \(10\)    & \checkmark \\
1000  & 0.1  & \(51\)    & \(-49\)    & \(-50\)   & \(50\)    & \checkmark \\
\bottomrule
\end{tabular}
\end{center}

Note that the sign flip is symmetric by construction, but the magnitude of the discrepancy grows linearly in \( L \), highlighting that the attribution gap between reference groups can be stretched arbitrarily far.

\subsection{Decision Tree Representation of \cref{eq:sign_flip_condition}}

  \begin{algorithm}[H]
    \caption{Unexplained component sign-flip}
    \label{alg:decision_tree}
    \begin{algorithmic}  
        \If{$\mu_H^T \Delta\beta \neq \mu_K^T \Delta\beta$}
            \If{$\sign(\mu_H^T \Delta\beta) = \sign(\mu_K^T \Delta\beta)$}
                \If{$\sign(\Delta\alpha) \neq \sign(\mu_H^T \Delta\beta)$}
                    \If{$|\mu_H^T \Delta\beta| < |\Delta\alpha| < |\mu_K^T \Delta\beta|$ \textbf{or} $|\mu_K^T \Delta\beta| < |\Delta\alpha| < |\mu_H^T \Delta\beta|$}
                        \State A sign flip occurs.
                    \Else
                    	\State There is no sign flip.
                    \EndIf
                \Else
                	\State There is no sign flip.
                \EndIf
            \ElsIf{$\sign(\mu_H^T \Delta\beta) \neq \sign(\mu_K^T \Delta\beta)$}
                \If{$\sign(\Delta\alpha) = \sign(\mu_H^T \Delta\beta)$ \textbf{and} $|\mu_K^T \Delta\beta| > |\Delta\alpha|$}
                    \State A sign flip occurs.
                \ElsIf{$\sign(\Delta\alpha) \neq \sign(\mu_H^T \Delta\beta)$ \textbf{and} $|\mu_H^T \Delta\beta| > |\Delta\alpha|$}
                    \State A sign flip occurs.
                \Else
                    \State There is no sign flip.
                \EndIf
            \EndIf
        \Else
            \State There is no sign flip.
        \EndIf
    \end{algorithmic}
\end{algorithm}


\subsection{Relationship Between Unexplained Component Sign Flips and Omitted Variables}
\label{sec:ovb}
Since unexplained component sign flips require \( | 1^T \Delta \beta | > | \Delta \alpha | \), any data-generating processes that jointly produces large \( | 1^T \Delta \beta | \) and \( \Delta \alpha \) would make sign flips less likely than \cref{prop:unexplained_fraction} suggests.

A candidate explanation is that the researcher's dataset omits important covariates, which are related to both \( X \) and \( Y \).
For example, hospital records only contain measurements that doctors have ordered, not the results of important tests a doctor may have skipped.
Similarly, census data used for studying wage gaps typically lack detailed information about skills and experience from a worker's resume and work history.
Finally, in practice, nonlinear transformations of the covariates are also effectively omitted from analysis, like when a researcher does not include polynomial terms in a regression specification.
Let \( Z \in \mathbb{R}^{d'} \) be a vector of such \textit{omitted covariates}.
We will establish that both \( \Delta \beta \) and \( \Delta \alpha \) can depend on the relationship between \( Z, X, \) and \( Y \).

Our argument requires us to specify the conditional expectations \( \E_g[Y \mid X, Z] \) and \( \E_g[Z \mid X] \).
To establish a simple example, we will assume that both are linear and satisfy
\begin{align}
  \label{eq:y_on_x_z}
  \E_g[Y \mid X, Z] &= \omega_g + X^\top \theta_g + Z^\top \gamma_g, \\
  \label{eq:z_on_x}
  \E_g[Z \mid X] &= \zeta_g + \Psi_g X.
\end{align}
Relaxing these assumptions may lead to other non-uniform relationships between \( \Delta \beta \) and \( \Delta \alpha \).

Recall that the OBD is based on regressing the outcome against the observed covariates for each group \(g \in \{H, K\}\) via
\begin{align}
    \label{eq:y_on_x}    
    \E_g[Y \mid X] = \alpha_g + X^\top \beta_g.
\end{align}
We will start by showing that \( \alpha_g \) and \( \beta_g \) depend on \( \gamma_g, \zeta_g, \), and \( \Psi_g \). 
To do so, apply the law of iterated expectations to \cref{eq:y_on_x}, plug in \cref{eq:y_on_x_z,eq:z_on_x}, and then rearrange terms:
\begin{align}    
    \E_g[Y \mid X] 
    &= \E_g[ \E_g[Y \mid X, Z] \mid X ] \nonumber \\
    &= \E_g[ \omega_g + X^\top \theta_g + Z^\top \gamma_g \mid X ] \nonumber  \\
    &= \omega_g + X^\top \theta_g + \E_g[Z \mid X]^\top \gamma_g \nonumber \\
    &= \omega_g + X^\top \theta_g + (\zeta_g + \Psi_g X)^\top \gamma_g \nonumber \\ 
    \label{eq:ovb}
    &= \underbrace{\omega_g + \zeta_g^\top \gamma_g}_{=\alpha_g} + \E_g[X]^\top \underbrace{(\theta_g + \Psi_g^\top \gamma_g)}_{=\beta_g}.
\end{align}
\cref{eq:ovb} is the well-known \textit{omitted variables bias} formula.
It describes how the slope and intercept of a regression equation depend on variables that are not included in the equation.
We can use this formula to express the difference in slopes and intercepts across the two populations as 
\begin{align}  
  \label{eq:delta_alpha_gamma}
  \Delta \alpha &= \Delta \omega + \Delta \zeta^\top \Delta \gamma,  \\
  \label{eq:delta_beta_gamma}
  \Delta \beta  &= \Delta \theta + \Delta \Psi^\top \Delta \gamma.
\end{align}
These equations show that \( \Delta \alpha \) and \( \Delta \beta \) depend on the role of omitted covariates in the data-generating process.
Crucially, they provide a mechanism to produce \( \Delta \alpha \) and \( 1^\top \Delta \beta \) that are jointly large.
If both the Z-X and Z-Y relationships change in complementary ways, then both \( \Delta \zeta^\top \Delta \gamma \) and \( 1^\top \Delta \Psi^\top \Delta \gamma \) will be large.


\subsection{Volume of the Sign-Flip Parameter Region}
We first show how to choose units for the covariates conveniently, so that \( \Delta \mu = 1 \).
This will simplify volume calculations for the regions of parameter space that flip the explained and unexplained components of the OBD.
\begin{lemma}
  \label{lemma:units}
  If \( \mu_H \neq \mu_K \) elementwise, we can choose units for \( X \) such that \( \mu_K := \E_K[X] = 0 \) and \( \mu_H := \E_H[X] = 1 \).
  \begin{proof}  
  Consider the transformation \( \tilde{X} = (X - \mu_K) / (\mu_H - \mu_K) \), where the division is elementwise.
  This is well-defined as long as \( \mu_H \neq \mu_K \) elementwise.
  Since \( \tilde{X} \) is a linear transformation of \( X \), it corresponds to a change in units.
  Under these new units, \( \E_K[\tilde{X}] = 0 \) and \( \E_H[\tilde{X}] = 1 \).  
  \end{proof}
\end{lemma}

\subsubsection{Proof of \cref{prop:explained_fraction}}
\label{proof:explained_fraction}
Apply \cref{lemma:units} so that \( \E_K[X] = 0 \), \( \E_H[X] = 1 \), and therefore \( \Delta \mu = 1 \).
\cref{eq:sign_flip_condition_explained_neq} now simplifies to 
\[ \sign\!\left(1^\top \beta_{H}\right) \neq \sign\!\left(1^\top \beta_{K}\right). \]
A Uniform$(-M, M)$ distribution on \( \beta_H \) and \( \beta_K \) corresponds to distributions on \( 1^\top \beta_{H} \) and \( 1^\top \beta_{K} \) that are identical and symmetric around zero.
As such, the signs of \( 1^\top \beta_{H} \) and \( 1^\top \beta_{H} \) will disagree in 50\% of the parameter space.

\subsubsection{Proof of \cref{prop:unexplained_fraction}}
\label{proof:unexplained_fraction}
Apply \cref{lemma:units} so that \( \E_K[X] = 0 \), \( \E_H[X] = 1 \), and therefore \( \Delta \mu = 1 \).
For the rescaled covariates, the conditions in \cref{thm:ordering_sensitivity} simplify.
In particular, the unexplained component flips sign if and only if \( 1^T \Delta \beta \neq 0 \), \( \sign (\Delta \alpha) \neq \sign (1^\top \Delta \beta) \), and \( | 1^T \Delta \beta | > |\Delta \alpha| \).

To compute the fraction of $C_M$ for which there are sign flips, we equivalently compute the probability of a sign flip under an independent $\mathrm{Uniform}(-M,M)$ measure on each element of the vector $(\beta_H, \beta_K, \alpha_H, \alpha_K) \in \R^{2d+2}$.
The condition \( 1^T \Delta \beta \neq 0 \) holds almost surely under the assumed Uniform$(-M,M)$ measure, since \( 1^\top \Delta \beta \) is a non-degenerate linear combination of continuously distributed random variables, so we can focus on the latter two conditions.
To do so, we will determine the distributions of  \( \Delta \alpha \) and \( 1^\top \Delta \beta \).

The sum of \( n \) Uniform$(0,1)$ random variables is sometimes called the Irwin-Hall(\(n\)) distribution.%
\footnote{For \(n = 2\) this is a familiar triangular distribution.}
When we place a Uniform$(-M, M)$ measure on the \( \alpha \)'s and \( \beta \)'s, there is a straightforward relationship between \( \Delta \alpha \), \( 1^\top \Delta \beta \), and appropriately defined Irwin-Hall random variables.
The difference \( \Delta \alpha \) is the sum of two Uniform$(-M, M)$ variables, which is equivalent to an Irwin-Hall($2$) distribution, rescaled to the interval $[-2M, 2M]$.
Similarly, \( 1^T \Delta \beta \) is the sum of \( 2 d \) Uniform$(-M, M)$ variables, and its distribution is an Irwin-Hall($2 d$) distribution, rescaled to the domain \( [-2dM, 2dM] \).

We can use this correspondence to translate the event 
\[
\{ |1^\top \Delta \beta| > |\Delta \alpha| \}
\cap
\{ \sign(\Delta \alpha) \neq \sign(1^\top \Delta \beta) \}.
\]
into an equivalent event in terms independent Irwin-Hall random variables.
Let \( I_2 \sim \text{Irwin-Hall}(2) \) and \( J_{2d} \sim \text{Irwin-Hall}(2d) \) be independent.
The above event is equivalent to 
\[
\bigl(\{ I_2 > 1 \} \cap \{ J_{2d} \le \tfrac{2dM - I_2}{2M} \}\bigr)
\;\cup\;
\bigl(\{ I_2 < 1 \} \cap \{ \tfrac{I_2 + 2dM}{2M} \le J_{2d} \}\bigr).
\]
That is, unexplained component sign flips have the same volume as the region where \( I_2 \) falls above its median and \( J_{2d} \) falls far in its left tail, or vice versa.
Because both \( I_2 \) and \( J_{2d} \) are symmetric, this volume is equal to twice the probability that \( I_2 \) is above its median and \( J_{2d} \) is sufficiently small.
That is:
\[
\Pr\!\left( \{ |1^\top \Delta \beta| > |\Delta \alpha| \} \right)
=
2\,\Pr\!\left( \{ I_2 > 1 \} \cap \left\{ J_{2d} < \tfrac{2dM - I_2}{2M} \right\} \right).
\]
We can calculate this exactly for small \( d \) by using the Irwin-Hall(n) cumulative distribution function (CDF),
\[ F(x) = \frac{1}{n!} \sum_{k=0}^{\lfloor x \rfloor} (-1)^k {n \choose k} (x - k)^n. \]
For larger \( d \), the combinatorial terms in the CDF become difficult to evaluate.
In those cases the central limit theorem allows us to approximate an Irwin-Hall(2d) distribution with a \( N(d, \frac{d}{6})\) distribution.
Finally, as \( d \to \infty \), sign flips have a volume equal to
\[ 2 Pr(I_2 > 1) Pr(J_{2d} < d \mid I_2 > 1).  \] 
Since the median of \( I_2 \) is $1$ and the median of \( J_{2d} \) is $d$, both probabilities are equal to \( \frac{1}{2} \), and the corresponding fraction of $C_M$ leading to sign flips is therefore \( 0.5 \).

\subsection{Proofs from \cref{sec:no_sign_flips}}


\subsubsection{Proof of \cref{prop:no_sign_flip}}
\label{proof:no_sign_flip}
Recall that by definition
$$
	\UK := \mu_H^T \Delta\beta + \Delta\alpha, \quad \quad \UH := \mu_K^T \Delta\beta + \Delta\alpha.
$$
By assumption, $\sign(\Delta\alpha) = \sign(\mu_H^T\Delta\beta) = \sign(\mu_K^T\Delta\beta)$, and thus $\sign(\UH) = \sign(\UK)$.
\subsubsection{Proof of \cref{prop:no_explained_sign_flip}}
\label{proof:no_explained_sign_flip}
As $\mu_K \neq \mu_H$ by assumption, \cref{lemma:units} applies to imply that we can fix units such that $\mu_K = 0$ and $\mu_H = 1$. Then $\Delta\mu = 1$, and the condition for the explained component to not sign flip, $\sign(\Delta\mu^T\beta_K) \neq \sign(\Delta\mu^T \beta_H)$ reduces to:
\begin{equation}
	\sign\left( \sum_{i=1}^d \beta_{K,i} \right) \neq \sign\left( \sum_{i=1}^d \beta_{H,i} \right),
\end{equation}
as claimed.

